\section{Atributos de Qualidade com GQM}\label{sec:gqm}
A Tabela~\ref{tab:gqm} aplica o GQM com quatro objetivos e três perguntas
por atributo. Para o NiA, trata-se de um plano
apoiado na inspeção; a coleta técnica foi aplicada ao aplicativo demonstrativo.

\begin{table*}[t]
\caption{Plano GQM. M: manutenibilidade; T: testabilidade; D: desempenho;
U: usabilidade. Medidas de tarefas com pessoas permanecem propostas.}
\label{tab:gqm}
\centering
\begin{tabularx}{\textwidth}{@{}lXXX@{}}
\toprule
Objetivo & Pergunta & Métrica/unidade & Estratégia de avaliação \\
\midrule
M: separar responsabilidades & Há ciclos entre módulos? & Quantidade de ciclos & Extrair e percorrer o grafo Gradle \\
M & Há dependências entre features? & Arestas feature--feature & Filtrar dependências diretas \\
M & A UI conhece a persistência? & Referências indevidas & Buscar tipos de Room/DAO nas features e UI \\
T: verificar com isolamento & Classes aceitam substitutos? & Classes testadas com fakes & Inspecionar casos e dependências dos testes \\
T & Os casos previstos passam? & Aprovados/executados & Consolidar relatórios JUnit \\
T & Quanto custa executar a suíte? & Tempo de parede (s) & Cronometrar a chamada de testes unitários \\
D: caracterizar execução & Quanto leva iniciar a Activity? & Mediana e intervalo (ms) & Dez reinícios do processo; \texttt{am start -W} \\
D & Quanto leva atualizar o fluxo? & Latência de emissão (ms) & Dez gravações alternadas com coletor ativo \\
D & Qual o uso após o roteiro? & PSS (MiB) & Uma captura \texttt{dumpsys meminfo} \\
U: avaliar clareza das ações & Controles têm rótulo e área? & Aprovados/5; mínimo 48 dp & Inspeção automatizada de cinco controles \\
U & Leitores concluem as tarefas? & Conclusões/tentativas (\%) & Sessões com pessoas: avaliação proposta \\
U & Quanto demoram nas tarefas? & Tempo por tarefa (s) & Cronometragem: avaliação proposta \\
\bottomrule
\end{tabularx}
\end{table*}

As métricas estruturais se relacionam à modularização e ao encapsulamento; as
de testes, às interfaces injetáveis; tempos e memória, à execução reativa; e
os controles, à apresentação. Nenhuma medida isolada comprova o atributo
inteiro. Sem comparação ou baseline, não se atribui causalidade à arquitetura.
